\section{Conclusion}

A role-owned choice is one event with several endpoint readings. Its owner
supplies a witness; the other endpoints consume that witness; all continue
with the same dependent instance. The public calculus makes this account
explicit through a global binder, local sum/product types, and a shared
substitution rule. Its execution, symmetry, and boundary-composition laws
follow from ordinary dependent computation.

Commitment adds a temporal distinction to that binding discipline. A value is
fixed under one audience and later made available to a larger one. The stable
binding connects the two events, while observation maps determine which
dependencies are permitted. Public factual outcomes become common knowledge
when the public readout is preserved by every participant's
indistinguishability relation.

The resulting foundation keeps the multi-role interaction intact. Objectives
can disagree, and analyses can select different notions of success, without
changing who owns a choice or which value a continuation received. The next
step is a dependent multiparty cut calculus whose open interfaces preserve
those same binding and observation laws.
